% This is text associated with the open textbook "Open Electromagnetics"
% See immediately below for author, date
% This work is licensed under the Creative Commons Attribution-ShareAlike 4.0 International License. (CC BY-SA 4.0) 
% To view a copy of the license, visit http://creativecommons.org/licenses/by-sa/4.0/

%ORB TITLE Convection and Conduction Currents
%ORB AUTHOR 0        # S.W. Ellingson, Virginia Tech (ellingson@vt.edu)
%ORB DATE 20180806   # yyyymmdd
%ORB ABSTRACT Current Distributions 

\label{m0110_Convection_and_Conduction} % <-- Must match name of this file and module directory name.  Don't mess with this!

In practice, we deal with two physical mechanisms for current: convection and conduction.  
The distinction between these types of current is important in electromagnetic analysis. 

{\bf Convection current} 
\index{current!convection}
consists of charged particles moving in response to mechanical forces, as opposed to being guided by the electric field (Sections~\ref{m0002_Electric_Field_Intensity} and/or \ref{m0102_Coulombs_Law}).  An example of a convection current is a cloud bearing free electrons that moves through the atmosphere driven by wind.

{\bf Conduction current} 
\index{current!conduction}
consists of charged particles moving in response to the electric field and not merely being carried by motion of the surrounding material.  
In some materials, the electric field is also able to dislodge weakly-bound electrons from atoms, which then subsequently travel some distance before reassociating with other atoms.  
For this reason, the individual electrons in a conduction current do not necessarily travel the full distance over which the current is perceived to exist. 

The distinction between convection and conduction is important because \emph{Ohm's Law} 
\index{Ohm's law}
(Section~\ref{m0010_Conductivity}) -- which specifies the relationship between electric field intensity and current -- applies only to conduction current.

%%%%%%%%%%%%%%%%%%%%%%%%%%%%%%%%%%%%%%%%%%%%%%%%%%%%%%%%%%%

{\bf Additional Reading:}
\begin{itemize}
\item ``\href{https://en.wikipedia.org/wiki/Electric_current}{Electric current}'' on Wikipedia.
\end{itemize}








